% Generated draft: requires mathematical and source review.
\documentclass[11pt]{article}
\usepackage[margin=25mm]{geometry}
\usepackage{fontspec}
\setmainfont{Latin Modern Roman}
\usepackage{amsmath,amssymb,mathtools,unicode-math}
\setmathfont{Latin Modern Math}
\usepackage{enumitem,needspace,xurl,hyperref}
\setlength{\parindent}{0pt}
\setlength{\parskip}{5pt}
\setlength{\emergencystretch}{3em}
\newcommand{\N}{\mathbb{N}}
\newcommand{\Z}{\mathbb{Z}}
\newcommand{\Q}{\mathbb{Q}}
\newcommand{\R}{\mathbb{R}}
\newcommand{\C}{\mathbb{C}}
\newcommand{\F}{\mathbb{F}}
\newcommand{\A}{\mathbb{A}}
\newcommand{\PP}{\mathbb{P}}
\newcommand{\E}{\mathbb{E}}
\newcommand{\eps}{\varepsilon}
\newcommand{\dd}{\,\mathrm{d}}
\begin{document}
\section*{UnsolvedMath ID: 30003285}
\textbf{Source problem code:} catalogue entry
\subsection{Definitions and mathematical statement}
Let \(G\) be a group with identity \(e\), and let \(\mu\) be a discrete probability distribution on \(G\). Symmetry means that \(\mu(g)=\mu(g^{-1})\) for every \(g\in G\). For independent increments \(Y_1,Y_2,\ldots\) with distribution \(\mu\), define \(X_0=e\) and \(X_n=Y_1\cdots Y_n\). Write \(\mu^{(n)}\) for the distribution of \(X_n\), so that \(\mu^{(2n)}(e)=\mathbb P(X_{2n}=e)\).

A bounded function \(f:G\to\R\) is \(\mu\)-harmonic if
\[
 f(x)=\sum_{g\in G}\mu(g)f(xg)\qquad(x\in G).
\]
The pair \((G,\mu)\) has the \emph{Liouville property} if every bounded \(\mu\)-harmonic function on \(G\) is constant.

\textbf{Return-probability question.} The supplied record asks whether a symmetric random walk satisfying a lower bound of the form
\[
 \mu^{(2n)}(e)\geq \exp(-c\sqrt n)
\]
necessarily has the Liouville property. Reading the bound as requiring a fixed constant \(c>0\), independent of \(n\), and validity for every positive integer \(n\), the proposed implication is
\[
 \left[\exists c>0\ \forall n\in\N,\ n\geq1:
 \mu^{(2n)}(e)\geq e^{-c\sqrt n}\right]
 \ \Longrightarrow\ 
 \left[\text{every bounded \(\mu\)-harmonic function on \(G\) is constant}\right].
\]
The source does not explicitly supply the quantifiers on \(c\) and \(n\). It also does not specify whether \(\operatorname{supp}\mu\) must generate \(G\), whether \(G\) must be countable or finitely generated, or whether \(\mu\) must satisfy finite-support, moment, or entropy assumptions. These are unresolved qualifications of the intended question; the displayed implication records a conventional reading of the bound, not a resolution of the missing hypotheses.

\textbf{Changing-generators question.} A conventional precise interpretation uses simple random walks on finite symmetric generating sets. For a nonempty finite set \(S\subseteq G\) with \(S=S^{-1}\) and \(\langle S\rangle=G\), put
\[
 \mu_S(g)=\begin{cases}|S|^{-1},&g\in S,\\0,&g\notin S.\end{cases}
\]
Under this interpretation, the question is whether, for every finitely generated group \(G\) and every two such sets \(S,T\),
\[
 (G,\mu_S)\text{ has the Liouville property}
 \quad\Longleftrightarrow\quad
 (G,\mu_T)\text{ has the Liouville property}.
\]
The supplied record does not specify symmetry of the generating sets, the rule assigning a step distribution to a generating set, or whether arbitrary symmetric weights are intended. Consequently this simple-random-walk formulation remains an explicitly qualified interpretation requiring source review.
\subsection{Short English statement}
The record asks whether a return probability bounded below by \(e^{-c\sqrt n}\) forces all bounded harmonic functions to be constant, and whether the Liouville property is invariant under changing a finite generating set. Essential assumptions on the first walk and the probability measures associated with generating sets remain unspecified. The supplied open status is dated curation, not independently verified here.
\subsection{Sources}
\begin{itemize}
\item Supplied local UnsolvedMath registry record, permanent ID \(30003285\), source problem code \texttt{catalogue entry}, entitled \emph{Return Probabilities and Stability of the Liouville Property}; collection: Oberwolfach Reports Open Problems. The record gives proposed year 2016 and identifies no individual proposer.
\item Supplied source citation: \emph{Heat Kernels, Stochastic Processes and Functional Inequalities} (2017), Oberwolfach Reports, pp.~3087--3147, DOI \texttt{10.4171/owr/2016/55}. \url{https://doi.org/10.4171/owr/2016/55}. The supplied background attributes the problem to workshop open-problem discussions documented in Oberwolfach Reports.
\item Supplied literature source: Peres and Zheng (2020). \url{https://doi.org/10.1093/imrn/rny034}. No article title or precise theorem statement is supplied.
\item Supplied literature source: Oberwolfach Report 55 (2016). \url{https://ems.press/content/serial-article-files/46659?nt=1}.
\item Literature assessment supplied with exact check date \texttt{2026-08-21}: the record labels the questions open and states that the sub-square-root return-probability threshold and stability under changing generators remain open, while later work proves only a strict sub-threshold criterion. This is supplied dated curation derived from the OWR material, not independently verified. The record itself requests independent verification of time-sensitive status and cited claims.
\end{itemize}
\end{document}
